% ============================================================
% CPEC2026 论文 LaTeX 模板（按《实验技术与管理》格式规范）
% ------------------------------------------------------------
% 说明：本模板为「非官方自建模板」。CPEC2026 征稿通知要求
%   “排版格式请参考附件投稿模板”，但该附件位于投审稿系统
%   (http://111.230.212.39/cpec/) 登录之后，公开站点不提供。
%   经核查，CPEC2026 与《实验技术与管理》合作、录用论文优先
%   发该刊，故本模板依据该刊官方投稿格式规范整理。
%   ⚠ 最终请以投审稿系统登录后的「附件投稿模板」为准。
% 编译：XeLaTeX（TeXstudio 默认编译器设为 XeLaTeX）
%   手动参考文献：单遍 XeLaTeX 即可；
%   改用 BibTeX（见文末注释）：XeLauncher→BibTeX→XeLaTeX×2
% ============================================================

\documentclass[12pt,a4paper]{ctexart}

% ---- 字体：正文宋体、标题黑体、西文/数字 Times New Roman ----
\usepackage{fontspec}
\usepackage{xeCJK}
\setmainfont{Times New Roman}
% 若系统装有 SimSun/SimHei，可显式指定（跨平台更稳）：
% \setCJKmainfont{SimSun}     % 宋体（正文）
% \setCJKsansfont{SimHei}     % 黑体（标题）

% ---- 页面（按《实验技术与管理》常见版式；以官方附件为准可调）----
% 上 3.5cm 给首页地脚留白；下 2.5cm；左右 2.8cm
\usepackage[top=3.5cm, bottom=2.5cm, left=2.8cm, right=2.8cm]{geometry}

% ---- 行距：正文约 20 磅（12pt×1.5≈18pt，取 1.5 近似）----
\usepackage{setspace}
\setstretch{1.5}

% ---- 中文段首缩进 ----
\usepackage{indentfirst}

% ---- 图表 ----
\usepackage{graphicx}
\usepackage{booktabs}     % 三线表：\toprule(顶) \midrule(栏目) \bottomrule(底)
\usepackage{array}
\usepackage{multirow}
\usepackage{caption}
\captionsetup{font=small, labelfont=bf, labelsep=space}
\usepackage{float}

% ---- 数学 ----
\usepackage{amsmath, amssymb}

% ---- 算法 / 代码（按需要保留或删除）----
\usepackage{algorithm}
\usepackage{algpseudocode}
\usepackage{listings}
\lstset{basicstyle=\ttfamily\small, breaklines=true, frame=single}

% ---- 矢量插图 ----
\usepackage{tikz}
\usetikzlibrary{positioning}
\usepackage{pgfplots}
\pgfplotsset{compat=1.18}

% ---- 超链接（黑白打印友好，关闭彩色）----
\usepackage[colorlinks=true, linkcolor=black, citecolor=black, urlcolor=black]{hyperref}

% ---- 标题层级：1 / 1.1 / 1.1.1（ctexart 默认即此），统一黑体 ----
\ctexset{
  section       = { format = \heiti\zihao{-3}\centering },
  subsection    = { format = \heiti\zihao{4} },
  subsubsection = { format = \heiti\zihao{-4} },
}

% ============================================================
% 以下为论文元信息：题目 / 作者 / 单位 / 首页地脚
% ============================================================
\title{AI赋能反诈实验教学案例设计与探索}

% 作者与单位；\textsuperscript{1} 对应下方单位编号，* 标记通信作者
% \thanks{...} 实现“首页地脚”（收稿日期/基金/作者简介/通信作者）
% 先清空脚注编号，避免标题页出现多余符号
\makeatletter
\renewcommand{\thefootnote}{}
\makeatother

\author{%
  韩冬\textsuperscript{1}，吴燕波\textsuperscript{1,*}，徐伟\textsuperscript{1}%
  \thanks{%
    收稿日期：2026-XX-XX\\
    基金项目：湖北省教育厅科研计划项目（项目编号待补）\\
    作者简介：韩冬（20XX—），男，湖北XX人，本科，研究方向为网络安全与反诈技术。E-mail：\href{mailto:xxx@xxx.edu.cn}{xxx@xxx.edu.cn}\\
    通信作者：吴燕波（19XX—），女，湖北XX人，教授，研究方向为计算机实践教育。E-mail：\href{mailto:xxx@xxx.edu.cn}{xxx@xxx.edu.cn}%
  }%
  \\%
  湖北警官学院~信息技术系，武汉~430000，中国%
}
\date{}

% ============================================================
\begin{document}

% 题名区上方留白微调（按需要）
\vspace*{-1.0cm}
\maketitle

% ------------------- 中文摘要 -------------------
\begin{abstract}
% 要求：第三人称；200–250 字；含“目的/方法/结果/结论”四要素；
% 禁用“本文/笔者”；不出现图、表、公式、参考文献。
针对公安院校计算机实践教学中反诈案例“难复现、难串并、难评价”的问题，
（目的）……（方法）……（结果）……（结论）……。

\par\noindent\textbf{关键词：}反诈教学；多智能体；图神经网络；实验案例设计；异构注意力
\par\noindent\textbf{中图分类号：}TP399 \quad \textbf{文献标识码：}A
\end{abstract}

% ------------------- 英文题名 / 作者 / 摘要 -------------------
\vspace{1em}
\begin{center}
  {\heiti\Large AI-Enabled Anti-Fraud Experimental Teaching Case Design and Exploration}

  \vspace{0.5em}
  {\normalsize HAN Dong\textsuperscript{1}, WU Yanbo\textsuperscript{1,*}, XU Wei\textsuperscript{1}}

  {\small Hubei University of Police, Department of Information Technology, Wuhan 430000, China}
\end{center}

\begin{center}\textbf{Abstract:}\end{center}
% 要求：与中文摘要对应；≥100 词；主动语态；避免第一人称；首句不与文题重复。
This paper addresses the difficulty of reproducing, correlating, and evaluating
anti-fraud cases in computer practice teaching at police colleges. ...
\begin{center}\textbf{Key words:}\end{center}
anti-fraud teaching; multi-agent; graph neural network; experimental case design; heterogeneous attention

% ============================================================
% 正文（层次标题用 1 / 1.1 / 1.1.1）
% ============================================================
\section{引言}
（此处撰写研究背景、教学痛点、本文贡献与结构。引言不编号或编为第 1 节均可。）

\section{相关工作}
\subsection{计算机实践实验教学研究}
（限三层标题：1.1 之下可用 1.1.1；再下用 1) 2) 或 (1)(2)。）

\section{实训载体：FraudLens 系统}
\subsection{系统总体架构}
（图随文给出，先见文字后见图。）

\begin{figure}[htbp]
  \centering
  % 图中文字用 Times New Roman 8pt；优先矢量图（pdf/emf/svg）
  \includegraphics[width=0.85\textwidth]{example-image.pdf}
  \caption{FraudLens 多阶段编排与质量门控降级架构}
  \label{fig:arch}
\end{figure}

\subsection{基于异构注意力的图神经网络}
（公式编号右对齐、全文连续；量符号首次出现时定义。）

\begin{equation}
  \beta_k = \mathrm{softmax}\bigl(q^{\top}\tanh(W h_k + b)\bigr),
  \label{eq:semantic}
\end{equation}

\section{实验与能力边界}
\subsection{基线对比}
（表格用三线表：顶线 1.5pt、栏目线 1pt、底线 1.5pt；数值按个位数对齐。）

\begin{table}[htbp]
  \centering
  \caption{不同方法在 Clean / Hard 场景下的 F1 对比}
  \label{tab:baselines}
  \begin{tabular}{lcc}
    \toprule[1.5pt]
    方法 & Clean F1 & Hard F1 \\
    \midrule
    KMeans        & 0.7821$\pm$0.0312 & 0.5431$\pm$0.0891 \\
    GraphSAGE    & 0.9345$\pm$0.0210 & 0.8762$\pm$0.0455 \\
    HAN（本文）   & 0.9488$\pm$0.0153 & 0.9154$\pm$0.1214 \\
    \bottomrule[1.5pt]
  \end{tabular}
\end{table}

\subsection{能力边界诚实标注}
（按贵刊规范，明确标注“未生效/场景依赖/单次运行不可靠”等边界，
 不把未验证机制当作已验证结论。）

\section{实训教学设计}
\subsection{学情分析与课程衔接}
（先修基础、分组角色、能力要求。）

\subsection{“研—训—评”模式与任务链}
（实验环节、课前-课中-课后、考核节点。）

\section{结语}
（总结贡献与局限，指出下一步工作。）

% ============================================================
% 参考文献（手动示例，展示《实验技术与管理》要求格式）
% 要求：顺序编码制；研究型论文≥12 条；
%       中文文献须附英文译文，并加注 (in Chinese)；
%       作者≤3 全列，>3 列前 3 加“, 等 / , et al”；
%       中文作者拼音：姓全大写、名缩写（例：HAN D Z）。
% ------------------------------------------------------------
% 【若改用你现有 BibTeX 工作流（gbt7714 + natbib + references.bib）】
%   1) 在导言区取消注释：\usepackage[numbers,sort&compress]{natbib}
%                          \usepackage{gbt7714}
%   2) 删除下方手动 thebibliography 块；
%   3) 此处改为：\bibliographystyle{gbt7714}
%                 \bibliography{../references}   % 草稿在 paper/ 下，模板在 templates_会议模板/
%   注意：gbt7714 默认不自动加 “(in Chinese)”，中文文献英文译文
%        需在 .bib 的 note 字段手动补全。
% ============================================================
\begin{thebibliography}{9}

\bibitem{liqingyong2024}
李清勇. “私教”还是“枪手”：基于大模型的计算机实践教学探索[J]. 实验技术与管理, 2024, 41(5): 1-8. (in Chinese)

\bibitem{shinn2023}
SHINN N, CASSANO F, LABBASH A, et al. Reflexion: language agents with verbal reinforcement learning[C]//Advances in Neural Information Processing Systems. 2023: 1-20.

\bibitem{wang2022han}
WANG X, JI H, SHI C, et al. Heterogeneous graph attention network[C]//The Web Conference. 2019: 2022-2032.

\bibitem{gbt7714}
中华人民共和国国家质量监督检验检疫总局, 中国国家标准化管理委员会. GB/T 7714-2015 信息与文献 参考文献著录规则[S]. 北京: 中国标准出版社, 2015. (in Chinese)

\end{thebibliography}

\end{document}
